% Part: first-order-logic
% Chapter: completeness
% Section: introduction

\documentclass[../../../include/open-logic-section]{subfiles}

\begin{document}

\iftag{FOL}
      {\olfileid{fol}{com}{int}}
      {\olfileid{pl}{com}{int}}

\olsection{परिचय}

संपूर्णता प्रमेय हा तर्कशास्त्राविषयीच्या सर्वांत मूलभूत निष्कर्षांपैकी एक
आहे. त्याच्या दोन मांडण्या आहेत; त्यांची सममूल्यता आपण सिद्ध करू. त्याची
पहिली मांडणी चिन्हार्थी परिणाम आणि आपल्या !!{derivation} पद्धतीतील संबंधाबद्दल
एक मूलभूत विधान करते: जर !!a{sentence}~$!A$ हे काही !!{sentence}s
$\Gamma$ यांच्यापासून निष्पन्न होत असेल, तर $\Gamma \Proves !A$ हे
प्रस्थापित करणारी !!a{derivation} सुद्धा असते. त्यामुळे प्रत्यक्षात निष्पन्न
न होणाऱ्या गोष्टी सिद्ध न करता !!{derivation} पद्धत जितकी सामर्थ्यवान असू
शकते तितकी ती असते.

दुसऱ्या मांडणीत ते प्रतिमानाच्या अस्तित्वाविषयीचा निष्कर्ष म्हणून सांगता येते:
!!{sentence}s चा प्रत्येक सुसंगत संच पूर्ततायोग्य असतो. सुसंगतता ही
सिद्धता-उपपत्तीय संकल्पना आहे: आपली !!{derivation} पद्धत विशिष्ट
!!{derivation}s निर्माण करू शकत नाही, असे ती सांगते. पण~$\Gamma$ पासून
विशिष्ट प्रकारच्या !!{derivation}s नाहीत एवढ्यावरूनच
\iftag{FOL}{!!a{structure}~$\Struct{M}$}{!!{valuation}~$\pAssign{v}$}
अस्तित्वात असून $\iftag{FOL}{\Sat{M}}{\pSat{v}}{\Gamma}$ आहे, याची हमी
आहे, असे कोण म्हणेल? संपूर्णता प्रमेय प्रथम सिद्ध होण्यापूर्वी---
खरे तर आज आपल्याकडे असलेल्या !!{derivation} पद्धती निर्माण होण्यापूर्वी---
महान जर्मन गणितज्ञ डाव्हीट हिल्बर्ट यांचे मत होते की गणिती सिद्धांतांची
सुसंगतता त्यांच्या विषयातील वस्तूंच्या अस्तित्वाची हमी देते. गोटलोप फ्रेग
यांना लिहिलेल्या पत्रात त्यांनी ते पुढीलप्रमाणे मांडले:
\begin{quote}
  स्वेच्छपणे दिलेली स्वयंसिद्धके आणि त्यांचे सर्व परिणाम परस्परविरोधी नसतील,
  तर ती सत्य असतात आणि स्वयंसिद्धकांनी परिभाषित केलेल्या वस्तू अस्तित्वात
  असतात. माझ्यासाठी सत्य आणि अस्तित्वाचा निकष हाच आहे.
\end{quote}
फ्रेग यांनी याला तीव्र विरोध केला. संपूर्णता प्रमेयाची दुसरी मांडणी दाखवते
की निदान पुढील अर्थाने हिल्बर्ट बरोबर होते: स्वयंसिद्धके सुसंगत असतील, तर
त्या सर्वांना सत्य करणारी \emph{एखादी}
\iftag{FOL}{!!{structure}}{!!{valuation}} अस्तित्वात असते.

संपूर्णता प्रमेय---किंवा अधिक नेमके सांगायचे तर त्याची सिद्धता---महत्त्वाची
असण्याची हीच काही एकमेव कारणे नाहीत. त्याचे अनेक महत्त्वाचे परिणाम आहेत;
त्यांपैकी काहींची स्वतंत्र चर्चा करू. उदाहरणार्थ, $\Gamma \Proves !A$ हे
दाखवणारी कोणतीही !!{derivation} सांत असल्यामुळे ती~$\Gamma$ मधील सांतच
!!{sentence}s वापरू शकते. म्हणून संपूर्णता प्रमेयावरून असे निष्पन्न होते की
$!A$ हा~$\Gamma$ चा परिणाम असेल, तर तो~$\Gamma$ च्या एखाद्या सांत उपसंचाचाही
परिणाम असतो. याला \emph{संहतता} म्हणतात. सममूल्य रीतीने, $\Gamma$ चा प्रत्येक
सांत उपसंच सुसंगत असेल, तर $\Gamma$ स्वतः सुसंगत असला पाहिजे.

संहतता प्रमेय हे संपूर्णता प्रमेयापासून !!{derivation}s मधून जाणाऱ्या वळणाने
मिळत असले, तरी संपूर्णता प्रमेयाची \emph{सिद्धता} वापरून ते थेट प्रस्थापित
करणेही शक्य आहे. कारण त्या सिद्धतेत विशिष्ट गुणधर्म---सुसंगतता---असलेला
!!{sentence}s चा संच घेऊन त्यातून विशिष्ट गुणधर्म असलेली !!a{structure}
रचली जाते (या प्रकरणात ती त्या संचाची पूर्ति करते). सुसंगत संचांऐवजी
``सांततः पूर्ततायोग्य'' !!{sentence}s च्या संचांपासून सुरुवात केल्यास जवळजवळ
हीच रचना वापरून संहतता थेट प्रस्थापित करता येते.
\iftag{FOL}{या रचनेतून इतर परिणामही मिळतात; उदाहरणार्थ, !!{sentence}s च्या
  कोणत्याही पूर्ततायोग्य संचाला सांत किंवा !!{denumerable}s असे प्रतिमान असते.
  (या निष्कर्षाला ल्योव्हेनहाइम--स्कोलेम प्रमेय म्हणतात.) सर्वसाधारणपणे,
  !!{sentence}s च्या संचांपासून !!{structure}s रचण्याची पद्धत तर्कशास्त्रात
  वारंवार आणि कधी कधी तत्त्वज्ञानातही वापरली जाते.}{}

\end{document}
